% In this file you should put the actual content of the blueprint.
% It will be used both by the web and the print version.
% It should *not* include the \begin{document}
%
% If you want to split the blueprint content into several files then
% the current file can be a simple sequence of \input. Otherwise It
% can start with a \section or \chapter for instance.


%%% GOALS %%%
% Big goal: prove F is finitely generated
% 1. Define F
% 2. Define elements of F
% 3. Tree pairs
% 4. Equivalence on binary trees (caret reduction)
% 5. Equivalence on tree pairs
% 6. Define generators
% 7. Infinite generating set
% 8. Reduction to finite.
% 9. Try doing both cantor set and tree pairs?
% 10. Equivalence trees and clopen subsets of Cantor
% 11. Equivalence tree pairs and cantor 
% 12. Do rational group instead?

\chapter{Preliminaries}

In this section we rehash some basics of the relations between prefix codes, trees and Cantor sets.

\section{Prefix codes}

Let $A$ be a set. Recall that a word over $A$ is a sequence of elements of $A$, the set of all words over $A$
including the empty word is denoted by $A^\ast$. We will write $xy$ for the concatenation of words $x, y\in A^\ast$. We denote the empty word by \(\varepsilon\).
A word $x$ is a \emph{prefix} of a word $y$ if there exists a word $z$ such that $xz = y$. We write $x \preceq y$
to denote that $x$ is a prefix of $y$.
\begin{definition}%
  \label{def:comparable}
  \rocq{Thompson.codes.comparable}
  \rocqok{}
  Let $x, y\in A^\ast$. We say that $x$ and $y$ are \emph{comparable} if $x$ is a prefix of $y$ or vice-versa, i.e.
  $x \preceq y$ or $y \preceq x$ holds. If $x$ and $y$ are not comparable, we say they are \emph{incomparable}.
\end{definition}

\begin{lemma}%
  \label{lem:comparable_symm}
  \uses{def:comparable}
  \rocq{Thompson.codes.comparable_symm}
  \rocqok{}
  Let $x, y\in A^\ast$. Then $x$ and $y$ are comparable if and only if $y$ and $x$ are comparable.
\end{lemma}
\begin{proof}
  \rocqok{}
  The proof follows immediately from the definition, as ``or'' is symmetric.
\end{proof}

An important role in the theory of Thompsons groups is played by prefix codes, which we now define.

\begin{definition}%
  \label{def:prefix_code}
  \uses{def:comparable}
  \rocq{Thompson.codes.prefix_code}
  \rocqok{}
  A prefix code over $A$ is a finite set of words $P$ over $A$ such that every pair of distinct
  words $w_1, w_2\in P$ are incomparable.
\end{definition}

\begin{definition}\label{def:concat_letter_set}
    Let \(x\in A\) and \(S\subseteq A^*\) be finite. We denote the set \(\makeset{xs}{\(s\in S\)}\) by \(xS\).
\end{definition}

\reinis{Finish this section}

\section{Trees}

\begin{definition}%
  \label{def:rooted_binary_tree}
\uses{def:prefix_code,def:concat_letter_set}
 A finite rooted binary tree \(T\) is one of:
 \begin{itemize}
     \item \(\{\varepsilon\}\);
     \item a set of the form \(0T_0 \cup 1T_1 \cup \{\varepsilon\}\). Where \(T_0\) and \(T_1\) are finite rooted binary trees.
 \end{itemize}
\end{definition}

\begin{definition}%
  \label{def:leaf_of_a_tree}
\uses{def:rooted_binary_tree}
If \(T\) is a finite rooted binary tree, then a leaf of \(T\) is one of
\begin{itemize}
    \item \(\varepsilon\) if \(T=\{\varepsilon\}\);
    \item \(0l\) if \(T=0T_0 \cup T_1 \cup \{\varepsilon\}\) and \(l\) is a leaf of \(T_0\);
    \item \(1l\) if \(T=0T_0 \cup T_1 \cup \{\varepsilon\}\) and \(l\) is a leaf of \(T_1\).
\end{itemize}
\end{definition}

\section{The Cantor set}

\chapter{Thompsons groups and monoids}

\chapter{Presentations}
